\section{Proof of Theorem~\ref{thm:ACP}}\label{sec:appendix-modules}
\begin{proof}[Proof of Theorem~\ref{thm:ACP}]
    The equivalence of (\ref{item:ofs-implies-Noetherian}) and (\ref{item:sub-ofs-is-ofs}) is a classical result in module theory. 
    Let us recall the proof. 
    For (\ref{item:ofs-implies-Noetherian}) $\Rightarrow$ (\ref{item:sub-ofs-is-ofs}), we build an orbit-finite spanning set for the subspace greedily, by adding orbit after orbit of vectors. 
    Thanks to (\ref{item:ofs-implies-Noetherian}), this must stabilise. 
    For (\ref{item:sub-ofs-is-ofs}) $\Rightarrow$ (\ref{item:ofs-implies-Noetherian}), we use (\ref{item:sub-ofs-is-ofs}) to get an orbit-finite spanning set for the union of the chain from (3), 
    and then we observe that orbits from this set can only be added finitely often in the chain. 

    Let us now prove the ``furthermore'' part. 
    \begin{itemize}
        \item 
        In one direction, consider equivariant subspaces $W \subseteq U \subseteq \Lin_\field \A^d$.
        As $\A$ is oligomorphic, the basis $\A^d$ is orbit-finite; 
        in particular $\Lin_\field \A^d$ is orbit-finitely spanned.
        But so is $U$ by~(\ref{item:sub-ofs-is-ofs}), as is the quotient $U/W$. 
        \item 
        In the other direction, take some orbit-finitely spanned vector space $V$. 
        By Definition~\ref{def:orbit-finite-set}, the spanning set can be obtained from some equivariant subset of $\A^d$ by quotienting under an equivariant equivalence relation. 
        This gives us a surjective equivariant linear map to $V$ from an equivariant subspace $U$ of $\Lin_\field \A^d$; 
        call its kernel $W$, so that we have a bijective equivariant linear map $V \leftarrow U/W$. 
    \end{itemize}

    Finally the implication (\ref{item:ofs-implies-Noetherian}) $\Rightarrow$ (\ref{item:LinAd-is-Noetherian}) is immediate, 
    whilst the converse implication follows from the ``furthermore'' part: 
    the lack of infinite chains is preserved by taking equivariant subspaces and images under equivariant linear maps. 
\end{proof}
